\documentclass[12pt]{article}
\usepackage[utf8]{inputenc}
\usepackage[T1]{fontenc}
\usepackage{hyperref}
\title{Krah\^o Language Resources for K-12 Curriculum Development}
\author{Compiled by Claude Sonnet 5.5 (claude-sonnet-5-5) for Andr\'as Kornai}
\date{}
\begin{document}
\maketitle

\section{Introduction}
Krah\^o (ISO 639-3 xra) is a Jê (Timbira) language of Brazil, spoken in Tocantins. It is still spoken as a first language in the community. My search was in Portuguese and was limited. It found material on Krah\^o and Timbira indigenous schooling, but no Krah\^o dictionary or grammar.

\section{Primary Educational Resources}
\begin{itemize}
\item ``Educa\c{c}\~ao Ind\'igena Krah\^o: material de apoio pedag\'ogico'' (article): \url{https://mail.sumarios.org/artigo/educa%C3%A7%C3%A3o-ind%C3%ADgena-krah%C3%B4-material-de-apoio-pedag%C3%B3gico}. According to the search summary, it describes bilingual, intercultural teaching materials produced at the Krah\^o school 19 de Abril in the village of Manoel Alves, in the Krah\^o language and Portuguese.
\item ``Escola Timbira'' (document on the Centro de Trabalho Indigenista library site), related to the Educa\c{c}\~ao Timbira collection of materials for Timbira villages including the Krah\^o. I did not open it, so I do not know its contents. \url{https://biblioteca.trabalhoindigenista.org.br/wp-content/uploads/2018/06/escola_timbira.pdf}
\end{itemize}

\section{Linguistic and Reference Materials}
I found no Krah\^o grammar or dictionary in this search. The search summary suggested that a glossary or dictionary may have been published by CTI or by the UFT Laborat\'orio de L\'inguas Ind\'igenas, but I did not confirm this.

\section{Religious and Community Resources}
None found.

\section{Digital and Audio Resources}
The search also returned articles on Krah\^o schooling. I did not read them, so I cannot say what they contain beyond their titles:
\begin{itemize}
\item \url{https://periodicos.ufpel.edu.br/index.php/rle/article/download/16066/10115}
\item \url{https://revistas.ufg.br/interacao/article/download/41003/21434/179638}
\item \url{https://revista.unitins.br/index.php/extensao/article/download/8369/4943}
\end{itemize}
I found no audio or software.

\section{Curriculum Development Context}
Krah\^o is a living language, and the school 19 de Abril already produces bilingual materials. A curriculum effort should start from that school and the Timbira education collection, working with Krah\^o teachers and the community. Studies cited in the search results point to Portuguese encroaching on the mother tongue in schools.
\end{document}